\section{Evaluation Design}\label{sec:method}
Our central claim is that \emph{consistency and correctness hold as the system scales across projects and agents}. We ask:
\begin{itemize}
\item \textbf{RQ1} Does the creation door keep the workspace consistent as the number of concurrently active projects grows from 1 to 10, compared with agents writing directly? (A scripted workload that \emph{demonstrates the mechanism}.)
\item \textbf{RQ2} With real LLM agents running concurrently, how often do they write outside their own project or folder, does the audit catch it, and do tasks succeed? How does this change after extending the door to the agents' own tools (Section~\ref{sec:p2})? (\emph{The main evidence}.)
\item \textbf{RQ3} Do the owner's own tests pass, before and after the changes? These are the lifecycle and adversarial tests and two LLM-driven tests.
\item \textbf{RQ4} How do the \hub's queue, callbacks and live sync behave: where does the time go, are callbacks delivered under failures, and does the mirror survive restarts?
\end{itemize}

\paragraph*{Setting and isolation} All experiments ran on a fresh, isolated copy of the production system with no credentials (all keys blank), every reference to the production folder rewritten, distinct ports, scheduled tasks disabled, and shims that block cloud and remote tools (Azure CLI, SSH, rsync, sshpass) for every process, including tools invoked by agents. The registry database is a private PostgreSQL cluster restored from a read-only dump of the production registry; the production database was only read. Every trial runs in a throw-away root with a freshly restored database. The machine is a MacBook Pro with an Apple M4 Pro (12 cores, 24\,GB RAM, macOS 26.5). The owner's production instance and its scheduled security agents kept running; we record their CPU load at the start and end of each experiment and report \emph{ratios} instead of absolute times, together with the CPU time and memory of the specific process tree under test.

\paragraph*{RQ1 workload} $N\in\{1,2,5,10\}$ projects, three agents each, all writing concurrently (12 writes per agent). The write mix follows the violation classes recorded in the system's history: 60\% valid, 15\% wrong sub-folder, 5\% project-root files, 10\% writes into \emph{another} project, 5\% duplicated utilities, 5\% references to another project's private file. Condition \textsc{door} sends every write through the real middleware; condition \textsc{bypass} writes to the requested path. We then run the real audit and heal. Five seeds per cell; 40 trials.

\paragraph*{RQ2 workload} $N\in\{1,2,5,10\}$ projects are created through the system's HTTP API, each with a Fullstack and a QA member, and all $N$ team orchestrations start at once. Each project gets a small, verifiable task: a Python module and a plain-assert test, with a module name unique to the project. From a content hash of the whole root before and after, we measure files changed inside the own project, inside another project, and outside any project. We also record the audit's findings, whether the module exists, and whether \emph{its test passes when we run it}. Cost and tokens come from the CLI's own accounting, and we record the share of wall time during which at least one model call was in flight. The model was fixed to a small model (\texttt{haiku}) for cost. Each $N$ ran three times \textbf{before} the changes of Section~\ref{sec:p2}. After the changes we ran the watch stage (one repetition per $N$ plus one extra at $N=10$) and then repeated the \textbf{same} protocol, three repetitions per $N$, with the hook enforcing (\textbf{after}). We also ran an \textbf{adversarial} variant at $N=5$: the task additionally asks each agent to copy its module into the shared library \emph{and} into another project's code folder. It ran once with the hook off and once with it on. For this variant we read each agent's own transcript to classify how a file reached shared code.

\paragraph*{RQ3} The owner's lifecycle test and adversarial test, unmodified, three repetitions each, before and after the changes. Two LLM-driven owner tests ran once before and once after: a multi-agent web build (``Revealer'') and a journey of three projects created from prompts. Cloud steps were run locally, and one prompt that asks for a remote sandbox deployment was left out.

\paragraph*{RQ4} Over real HTTP against the \hub's stand-in server: third-party clients submit $M$ requests (50 with one client, 200 with four), a poller pulls, processes (synthetic, zero time, to isolate overhead) and publishes; a local receiver acknowledges callbacks, optionally failing 20\% of them. Each request is time-tagged at \emph{in, queued, pulled, processed, published} and \emph{callback}. For live sync we push 20 pool updates and time push-to-visible, then restart the \hub{} and check the mirror. Three repetitions.

\paragraph*{Regression suite} A suite of 18 checks was baselined before any change. It covers isolation invariants, compilation, live endpoints returning real data, the door's placement and owner-scope rules, the audit's detection classes, router determinism, the queue contract, and sequence contiguity. With the changes it grew to 32 checks: hook decisions (allow, redirect, block, door use and actor check, and two parser false positives found during the study), clean project deletion, and a resources section for every role. It was re-run after every change to the isolated copy. The final run passed 32/32 with no regression against the baseline.
